% College Lecture 13: ElGamal and ECC. All points and values verified by script.
\section{ElGamal and Elliptic Curve Cryptosystems}

\subsection{ElGamal}
Security rests on the \textbf{discrete logarithm problem}: $e_2 = e_1^d \bmod p$ is easy; finding $d = \log_{e_1} e_2 \bmod p$ is infeasible for large $p$.
\begin{keybox}[title=ElGamal]
\textbf{Keys}: large prime $p$; $e_1$ a primitive root mod $p$; private $d$ with $1 \le d \le p - 2$; \;$e_2 = e_1^d \bmod p$. Public $(e_1, e_2, p)$, private $d$.\\
\textbf{Encrypt} (fresh random $r$): $C_1 = e_1^r \bmod p$, \;$C_2 = P \cdot e_2^r \bmod p$.\\
\textbf{Decrypt}: $P = C_2 \cdot (C_1^d)^{-1} \bmod p$. \;Proof: $C_1^d = e_1^{rd} = e_2^r$, so $C_2 (C_1^d)^{-1} = P e_2^r (e_2^r)^{-1} = P$.
\end{keybox}
\textbf{Example}: $p = 11$, $e_1 = 2$, $d = 3$ $\Rightarrow e_2 = 8$. Alice sends $P = 7$ with $r = 4$: $C_1 = 2^4 \bmod 11 = 5$, \;$C_2 = 7 \cdot 8^4 \bmod 11 = 7 \cdot 4096 \bmod 11 = 6$. Bob: $C_1^d = 5^3 \bmod 11 = 4$, $4^{-1} = 3$, $P = 6 \cdot 3 \bmod 11 = \mathbf{7}$.\\
\textbf{Shortcut}: $(C_1^d)^{-1} = C_1^{p-1-d}$ (Fermat): $5^{7} \bmod 11 = 3$ \checkmark.

\textbf{Security}: \textbf{low-modulus attack} (small $p$ lets Eve solve the discrete log once and decrypt everything; use $p \ge 1024$ bits); \textbf{known-plaintext attack} if $r$ is reused: from $C_2 = P e_2^r$ she gets $e_2^r = C_2 P^{-1}$, then $P' = C_2' (e_2^r)^{-1}$. Use a fresh $r$ every time. Note the ciphertext is twice the size of the plaintext.

\subsection{Elliptic curves over $\mathbb{Z}_p$}
RSA and ElGamal need big keys; ECC gives the same security with much smaller keys (160-bit ECC $\approx$ 1024-bit RSA).
\[
E_p(a, b):\ y^2 \equiv x^3 + ax + b \pmod p,\quad p > 3,\quad 4a^3 + 27b^2 \not\equiv 0 \pmod p,\ \text{ plus the point at infinity } O.
\]
\textbf{Finding points}: for each $x$, compute $w = x^3 + ax + b \bmod p$; if $w$ is a square mod $p$, the points are $(x, \pm\sqrt w)$. $E_{13}(1,1)$:
\begin{center}\small
\ttx{(0,1) (0,12) (1,4) (1,9) (4,2) (4,11) (5,1) (5,12) (7,0) (8,1) (8,12) (10,6) (10,7) (11,2) (11,11) (12,5) (12,8)}
\end{center}
17 points $+ O = 18$. $x = 2, 3, 6, 9$ give non-squares. Inverse of $(x, y)$ is $(x, -y) = (x, p - y)$; $(7, 0)$ is its own inverse.

\begin{keybox}[title=Point addition and doubling]
$R = P + Q$: \quad $x_3 = s^2 - x_1 - x_2$, \quad $y_3 = s(x_1 - x_3) - y_1 \pmod p$, where
\[
s = \frac{y_2 - y_1}{x_2 - x_1}\ (P \ne Q,\ \text{addition}),\qquad s = \frac{3x_1^2 + a}{2y_1}\ (P = Q,\ \text{doubling}).
\]
Division means multiplying by the inverse mod $p$. $P + (-P) = O$; $P + O = P$.
\end{keybox}
\textbf{Doubling on $E_{17}(2, 2)$}, $P = (5, 1)$: $s = (3 \cdot 25 + 2) \cdot (2)^{-1} = 77 \cdot 9 \equiv 9 \cdot 9 = 81 \equiv 13$; \;$x_3 = 169 - 10 = 159 \equiv 6$; \;$y_3 = 13(5 - 6) - 1 = -14 \equiv 3$ $\Rightarrow 2P = (6, 3)$. (This $P$ generates all 19 points of the curve.)

\textbf{Scalar multiplication} $kP = P + \dots + P$: in $E_{13}(1,1)$, $4 \cdot (1, 4) = (11, 11)$ and $3 \cdot (8, 1) = (10, 7)$.

\textbf{Hasse}: $p + 1 - 2\sqrt p \le \#E \le p + 1 + 2\sqrt p$. \textbf{ECDLP}: given $P$ and $T = dP$, find $d$; best known method is Pollard rho, infeasible for large curves.

\subsection{ECC simulating ElGamal}
\begin{keybox}
\textbf{Keys}: curve $E$, point $e_1$, private $d$, \;$e_2 = d \times e_1$. Public $(E, e_1, e_2)$.\\
\textbf{Encrypt} (plaintext is a point $P$, random $r$): $C_1 = r \times e_1$, \;$C_2 = P + r \times e_2$.\\
\textbf{Decrypt}: $P = C_2 - d \times C_1$ \;(subtract $=$ add the inverse point), since $d \times r \times e_1 = r \times e_2$.
\end{keybox}
\textbf{Example} $E_{67}(2, 3)$: $e_1 = (2, 22)$, $d = 4$ $\Rightarrow e_2 = (13, 45)$. Alice: $P = (24, 26)$, $r = 2$: $C_1 = (35, 1)$, $C_2 = (21, 44)$. Bob: $4 \times C_1 = (23, 25)$, inverse $(23, 42)$, $(21, 44) + (23, 42) = \mathbf{(24, 26)}$ \checkmark.

\begin{trap}[: point arithmetic slips]
The doubling slope has $3x_1^2 + a$ (uses $a$, never $b$). Compute $s$ mod $p$ with a modular inverse, never a fraction. $-P$ flips only $y$.
\end{trap}
